\newcommand{\verfasser}{Phil Steinhorst}					% Verfasser des Dokuments
\newcommand{\fach}{Höhere Algebra II}				% Titel der Vorlesung								
\newcommand{\shortFach}{HAlg2}								% Titel kurz
\newcommand{\prof}{Prof.\,Dr.\,Lutz Hille}					% Dozent
\newcommand{\untertitel}{gelesen von \prof}					% Untertitel
\newcommand{\semester}{Wintersemester 2014/2015}					% Semester
\newcommand{\homepage}{}	% Vorlesungshomepage
% % % % % % % % % % % % % % % % % % % % % % % % % % % % % % % % %
%		PhistScript.tex											%
%		XeLaTeX-Settings for Scripts							%
%																%
%		Author: Phil Steinhorst									%
% % % % % % % % % % % % % % % % % % % % % % % % % % % % % % % % %

\setlength\parindent{0pt}             % ohne Einrueckung

% % % Konfiguration von scrheadings
\setheadsepline{1pt}[\color{black}]
\pagestyle{scrheadings}
\clearscrheadfoot

% % % Kopf-/Fußzeilenlayout
\providecommand{\shortFach}{\fach}
\ohead{\small \leftmark}
\ofoot[{ \small \thepage}]{ \small \thepage} 
\automark{section}

% Metadaten
\title{\fach}
\author{\verfasser}
\subtitle{\untertitel}
\date{\semester}

% Inhaltsverzeichnis
\setcounter{tocdepth}{1}
\usepackage[tocindentauto]{tocstyle}
\usetocstyle{KOMAlike}	

% Nummerierung
\numberwithin{equation}{section}		% Section in Equation-Nummerierung einbeziehen

% ntheorem package
\newcounter{tcount}						% Zähler für Theoreme
\numberwithin{tcount}{section}			% Section in Theorem-Nummerierung einbeziehen
\usepackage[hyperref]{ntheorem}			% Lade ntheorem mit hyperref-Workaround
\theoremstyle{break}					% Zeilenumbruch nach Theoremkopf
\theorembodyfont{\normalfont}			% nicht kursiv
\theorempreskipamount0.5cm				% Abstand vor Theorem
\theorempostskipamount0.5cm				% Abstand nach Theorem
\newtheorem{thm}[tcount]{Theorem}
\newtheorem{bsp}[tcount]{Beispiel}
\newtheorem{defn}[tcount]{Definition}
\newtheorem{bem}[tcount]{Bemerkung}
\newtheorem{lemma}[tcount]{Lemma}

% Workaround für Tabulatoren im Theorem-Verzeichnis
\makeatletter
\def\thm@@thmline@name#1#2#3#4{%
        \@dottedtocline{-2}{0em}{2.3em}%
                   {\makebox[\widesttheorem][l]{#1 \protect\numberline{#2}}#3}%
                   {#4}}
\@ifpackageloaded{hyperref}{
\def\thm@@thmline@name#1#2#3#4#5{%
    \ifx\\#5\\%
        \@dottedtocline{-2}{0em}{2.3em}%
            {\makebox[\widesttheorem][l]{#1 \protect\numberline{#2}}#3}%
            {#4}
    \else
        \ifHy@linktocpage\relax\relax
            \@dottedtocline{-2}{0em}{2.3em}%
                {\makebox[\widesttheorem][l]{#1 \protect\numberline{#2}}#3}%
                {\hyper@linkstart{link}{#5}{#4}\hyper@linkend}%
        \else
            \@dottedtocline{-2}{0em}{2.3em}%
                {\hyper@linkstart{link}{#5}%
                  {\makebox[\widesttheorem][l]{#1 \protect\numberline{#2}}#3}\hyper@linkend}%
                    {#4}%
        \fi
    \fi}
}
\makeatother
\newlength\widesttheorem
\AtBeginDocument{
  \settowidth{\widesttheorem}{Proposition 10.10\quad}	% Referenzstring für die Breite der ersten Spalte
}

\theoremlisttype{optname}	% nur Theoreme auflisten, die einen Namen haben. 							% XeLaTeX-Konfiguration für Skripte

% % % Dokumentspezifische Einstellungen % % %
\setcounter{tocdepth}{2}				% Tiefe im Inhaltsverzeichnis (1: nur Sections, 2: auch subsections...)
\numberwithin{equation}{section}		% Section in Equation-Nummerierung einbeziehen
\newcounter{countsatz}						% Zähler für Sätze
\numberwithin{countsatz}{section}			% Section in Satz-Nummerierung einbeziehen
\newcounter{countfalko}						% Zähler für Fs
\numberwithin{countfalko}{section}			% Section in F-Nummerierung einbeziehen
\newcounter{countbsp}						% Zähler für Beispiele
\numberwithin{countbsp}{section}			% Section in Beispiel-Nummerierung einbeziehen
\newcounter{countdef}						% Zähler für Definitionen
\numberwithin{countdef}{section}			% Section in Definition-Nummerierung einbeziehen
\newcounter{countlemma}						% Zähler für Lemmata
\numberwithin{countlemma}{section}			% Section in Lemma-Nummerierung einbeziehen

% ab hier können neue Umgebungen für Sätze definiert und benannt werden!
% Syntax: \newtheorem{interner Umgebungsname}[Zähler]{Name im Dokument}
% Möchte man Sätze und Definitionen getrennt durchnummerieren, benötigt man weitere Zähler.
\newtheorem{satz}[countsatz]{Satz}
\newtheorem{bsp}[countbsp]{Beispiel}
\newtheorem{defn}[countdef]{Definition}
\newtheorem{lemma}[countlemma]{Lemma}
% % % % % % % % % % % % % % % % % % % % % % %

\begin{document}
% % % % % % % % % % % % % % % % % % % % % % % % % % % % % % % % %
%		PhistTitle.tex											%
%		Layout für Skript-Titelseiten							%
%																%
%		Author: Phil Steinhorst									%
% % % % % % % % % % % % % % % % % % % % % % % % % % % % % % % % %
\begin{titlepage}
\pagestyle{empty}
\begin{center}
\begin{minipage}{0.4\textwidth}
\begin{flushleft}
\includegraphics[height=1.5cm,keepaspectratio]{../!config/wwulogo.pdf}
\end{flushleft}
\end{minipage}
\hfill
\begin{minipage}{0.4\textwidth}
\begin{flushright}
\vspace*{0.3cm}
\includegraphics[height=1.2cm,keepaspectratio]{../!config/fb10logo.pdf} \
\end{flushright}
\end{minipage}

% Title
\vspace*{6cm}
\textbf{\Huge{\fach}} \\
\vspace{0.2cm} 
\textbf{{\LARGE \untertitel}} \\
\vspace{0.6cm}
\LARGE{Mitschrift von \verfasser} \\
\vspace{0.6cm}
\LARGE{\semester} \\
\vspace*{2cm}
\footnotesize{\url{\homepage}}
\vfill

\vspace*{2cm}
\begin{flushright}
{\footnotesize Stand: \today}
\end{flushright}
\end{center}
\end{titlepage}							% Titelseite ausgeben
\section*{Vorwort}
\label{sec:preface}
	Der vorliegende Text ist eine inhaltliche Aufbereitung zur Vorlesung Elliptische Kurven und Kryptographie, gelesen von PD Dr. Karin Halupczok an der WWU Münster im Sommersemester 2015. Der Inhalt entspricht weitestgehend dem handschriftlichen Skript, welches auf der Vorlesungswebsite bereitsgestellt wird. Dieses Werk ist daher keine Eigenleistung des Autors und wird nicht von der Dozentin der Veranstaltung korrekturgelesen. Für die Korrektheit des Inhalts wird keinerlei Garantie übernommen. Bemerkungen, Korrekturen und Ergänzungen kann man folgenderweise loswerden:
	\begin{itemize}
		\item persönlich durch Überreichen von Notizen oder per E-Mail
		\item durch Abändern der entsprechenden TeX-Dateien und Versand per E-Mail an mich
		\item direktes Mitarbeiten via GitHub. Dieses Skript befindet sich im \texttt{latex-wwu}-Repository von Jannes Bantje:
		\begin{center}
			\url{https://github.com/JaMeZ-B/latex-wwu}
		\end{center}
	\end{itemize}

\subsection*{Literatur}
\label{sub:lit}
\begin{itemize}
	\item Blake, Seroussi, Smart: Elliptic curves in cryptography
	\item Menezes, van Oorschot, Vanstone: Handbook of applied cryptography
	\item Silverman: The arithmetic of elliptic curves
	\item Silverman: A friendly introduction to number theory, chap. 40-45
	\item Washington: Elliptic curves, number theory and cryptography
	\item Werner: Elliptische Kurven in der Kryptographie
\end{itemize}

\subsection*{Kommentar der Dozentin}
In der Vorlesung beschäftigen wir uns mit den arithmetischen und geometrischen Eigenschaften elliptischer Kurven sowie deren Anwendungen in der Kryptographie. Dabei werden wir auch einen Vergleich mit Anwendungen der elementaren Zahlentheorie in der Kryptographie ziehen. Wir verfolgen eine elementare Herangehensweise, d.h. Kenntnisse der algebraischen Geometrie und der Funktionen- oder Zahlentheorie werden nicht benötigt. Es genügen die Vorkenntnisse aus den Grundvorlesungen.

\subsection*{Vorlesungswebsite}
\label{sub:link}
Das handgeschriebene Skript sowie weiteres Material findet man unter folgendem Link:
\begin{center}
	\url{\homepage}
\end{center}

\subsection*{Titelbild}
\label{sub:titlepic}
Das fehlt noch. Über Ideen und Anregungen freue ich mich sehr!

\vfill
\begin{flushright}
	Phil Steinhorst \\
	p.st@wwu.de
\end{flushright}
\newpage
\tableofcontents
\newpage											% Inhaltsverzeichnis ausgeben
\chapter{CAT(0)-Räume} % (fold)
\label{cha:1}

\begin{definition}[metrischer Raum]
	Sei $X$ eine Menge. Eine Abbildung $d \colon X \times X \rightarrow \RR_{\geq 0}$ heißt \Index{Metrik}, wenn für alle $x,y,z \in X$ gilt:
	\begin{enumerate}[(i)]
		\item $d(x,y) = 0 \Leftrightarrow x=y$
		\item $d(x,y) = d(y,x)$
		\item $d(x,z) \leq d(x,y) + d(y,z)$
	\end{enumerate}
	Das Paar $(X,d)$ heißt dann \Index{metrischer Raum}.
\end{definition}

\cleardoubleoddemptypage

%\appendix							% Anhang beginnt hier (Nummerierung mit Buchstaben)
\printindex							% Index ausgeben
\subsection*{Liste der Sätze und Definitionen}
\listtheorems{satz,defn,lemma}					% listet angegebene Sätze mit Namen und Seitenzahl auf
\end{document}